% !TeX root = ../papers/nc-p3.tex
\section{可除指数による四次曲面像と滑らかな二次三次元体上の基点の排除}
\label{sec:exclude-del-pezzo-center-images}
\begin{theorem}\label{thm:exclude-del-pezzo-center-images}
定理\ref{thm:quartic-center-image-atlas}の第三型は正則整域窓に現れない。
第二型で $Q$ が滑らかなら常に四次系は基点無しである。
従って四次中心系に基点を持つ十五次元閉包の正則整域窓の像として残るのは、
特異な三次元二次曲面への双有理写像と三次スクロール曲面への写像だけである。
\end{theorem}
\begin{proof}[次数四デルペッツォ曲面上の可除指数]
第三型の完全交叉曲面を $Y$、斉次座標環を $P$ とする。二次の補元を選んで平方完成すると
$Z$ は $P[u]$ のモニックな二重被覆で、有理特異点を持つ $Z$ の対合不変環が $P[u]$ である。
有限群商で有理特異点は保たれるので $P$ も有理特異点を持つ。
従って $Y$ はGorenstein有理特異点、すなわちDu Val点だけを持つ次数4デルペッツォ曲面である。
最小解消は滑らかな有理弱デルペッツォ曲面である。$U=Y_{\rm sm}$ とすれば
$\operatorname{Br}(U)=0$。実際有理な解消での剰余を調べると例外は互いに交わらない
有理曲線の木であり、葉から順に分岐が消えて全ての剰余が零になる。解消自体のBrauer群は零である。

$U$ の有限巡回指標は全て位数2以下であることも記す。単元は定数で、
$\operatorname{Pic}(U)$ は解消のPicard群を例外曲線の根格子で割ったものである。
弱デルペッツォ曲面の吹上げ格子で標準類の直交は $D_5$ であり、各例外根は
$\pm e_i\pm e_j$ ($i\ne j$) となる。例えば平面の五点の吹上げ基底で
連続した例外差四本と直線から最初の三例外を引いた根が $D_5$ の単純根を与える。
生成根を符号付き辺とし、各連結成分の全域木で座標の符号を揃えて差の根とする。
残る辺も差なら成分の格子は座標和零の全格子で飽和している。
和の辺が一つでもあれば差と合わせ座標和偶数の全格子を生成する。
従って飽和との商は2倍で零である。飽和は標準類と直交するので、この計算で全ての捩れを覆う。
Kummer列が指標の主張を与える。

$u,w$ の最高二次形式が相異なる二根を持てば、二変数の線型変更および $x$ の二次式だけの平行移動により
$uv=f_4(x)$ とできる。射影中心の $U$ 上で $u\ne0$ とした部分は直線束
$\mathcal O_U(2)$ の零切断を除いた部分である。そこでは滑らかで次数1の局所単元があるから、
指数4の一般Brauer類 $\alpha$ はこの部分で非分岐である。
従って零切断での剰余 $\chi$ は $U$ 全体の非分岐巡回指標であり、位数は高々2である
（直線束の全空間での余次元二剰余を使う）。$\chi$ の与える連結エタール被覆 $U'\to U$ に引き戻すと、
零切断へも純性で延びる。直線束全体のBrauer群は底のBrauer群に等しい。
$U'$ のBrauer群も零である。実際 $Y$ のその体での正規化は余次元1でエタールな次数高々2の被覆で、
Du Val点の穴あき局所被覆は有限 $\mathrm{SL}_2$ 群の部分群による同じ型の商の穴あき部分である。
よって正規化は再びDu Valで、反標準因子は引き戻しで豊富、その解消は有理弱デルペッツォ曲面である。
有限個の滑らかな点を除いてもBrauer群は変わらず、同じ木の議論を使える。
従って $\alpha$ は次数高々2の拡大で分裂し、指数4に反する。

二重根の場合は $w^2=u\ell_2(x)+f_4(x)$ と書ける。$\ell_2=0$ なら $Z$ は
一般八次曲面の標準環の多項式拡大になる。この標準環の頂点は有理特異点でない。
実際標準束は豊富で、解消上に引き戻した錐の零切断から
$H^2(S,\mathcal O_S)\ne0$ が出る（同じく有理特異点の切断環の正次数・零次数コホモロジー消滅を使える）。
これは $Z$ の有理特異点性に反する。$\ell_2\ne0$ のとき、その非消滅開集合上では $w$ は直線束の
自由座標であり、$\alpha$ は $\beta\in\operatorname{Br}(k(Y))$ の引き戻しである。
一般繊維のアフィン直線全体で非分岐なのでアフィン直線のBrauer不変性による。
$u$ を自由座標にした別の直線束の上へは $w$ による有限二重被覆である。
$Y$ の任意の素曲線の一般点で底を滑らかに取り、この直線束内の垂直因子を考える。
その上の全ての因子で $\alpha$ は非分岐である。被覆は正規、余次元1で滑らかで、
$x$ のいずれかが非零だから次数安定化もない。剰余とcorestrictionの両立から
下にある $2\beta$ の剰余は零となる。垂直因子の体は曲線の関数体の純超越拡大なので
曲線上の剰余自身が零である。従って $2\beta\in\operatorname{Br}(U)=0$。
代数閉体上の曲面関数体におけるde Jongの周期・指数定理により $\beta$ の指数は高々2で、再び矛盾する。
\end{proof}

\begin{proof}[滑らかな三次元二次曲面上の全ての極因子]
第二型で基点があり $Q$ が滑らかとする。$\alpha\in\operatorname{Br}(k(Q))[4]$ は
$T=(l_2=0)$ の外で非分岐である。$T$ が滑らかな場合は前節で除いた。
$T$ が被約既約で特異な場合、特異点 $p$ で $l_2$ から $q_2$ の定数倍を引き、$p$ を根基とする二次式を得る
（根基がより大きくてもよい）。$p$ の接超平面を無限遠に選ぶと
$Q$ の残りは $\mathbb A^3$ である。この射影座標では $p$ の成分だけが残り三座標の二次式で表され、
新しい極方程式はその成分を使わない。
従ってこのアフィン三空間上の唯一の分岐候補は、射影閉包が $\mathbb P^3$ の階数4または3の二次曲面である。
階数2以下なら元の $T$ が可約または重複になる。
Brauer剰余はその滑らかな部分に非分岐な巡回指標であり、その指標は $\operatorname{Br}(\mathbb A^3)=0$ により
$\alpha$ を一意に決める。

階数4では射影二次曲面は $\mathbb P^1\times\mathbb P^1$、無限遠切断は滑らかな $(1,1)$ または
二罫線の和である。どちらでもそのアフィン部分の単元は定数でPicard群は捩れなしである。
実際境界の成分の類は前者なら原始、後者なら基底である。従って指標はない。
階数3で頂点がアフィン部分にあるとき、座標を変えて方程式は $F=xy-z^2$ となる。
穴あき二次錐の非自明な巡回指標は一つの位数2のものだけである
（$\mathbb A^2\setminus\{0\}$ を符号で割った普遍二重被覆）。その指標は $\sqrt{x}$ で表せる。
四元数記号 $(F,x)_2$ の剰余がちょうどこれで、$x=0$ では $F=-z^2$ が平方なので他の分岐はない。
よって $\alpha$ の指数は高々2である。

階数3の頂点が無限遠なら、アフィン極因子は滑らかなアフィン二次曲線とアフィン直線の積である。
二次曲線に無限遠が一点なら指標はない。二点ならアフィン線型座標で $F=XY-1$ とできる。
指標は $X$ の冪の四乗根で表されるから、$\alpha$ は記号 $(F,X)_4^j$ に等しい。
実際 $X=0$ での剰余は定数 $-1$ の類で零であり、残りの全剰余が同じである。
元の $Q$ の接超平面因子は極因子の成分でなかったので、そこでの剰余も零でなければならない。
$X,Y$ はそれぞれ一次の極を持ち、$F$ は二次の極を持つため、剰余は
$(X/Y)^j$ の先頭係数の類（符号・逆元は無関係）である。
接超平面切断は $p$ を頂点とする滑らかな平面二次曲線上の錐である。
二つの一次主要部の比はこの二次曲線上の独立な一次式の比で、約分後の次数は1または2である。
独立性は $X,Y$ のアフィン線型部の独立性と平面二次曲線の非退化性による。
この非定数有理関数は四乗でなく、錐の純超越射影方向を加えても四乗にならない。
ゆえに $j$ は偶数で、記号は自明または四元数に等しい。指数4は不可能である。

残る重複または可約な $T$ は、$\operatorname{Pic}(Q)=\mathbb Z[H]$ より
一つの超平面の二倍または二超平面の和である。一成分だけならその滑らかな部分に非自明指標はない。
滑らかな二次曲面では明らかで、射影二次錐の頂点を除いた部分も $\mathbb P^1$ 上の直線束の全空間である。
二成分の一方が滑らかな場合、その成分での剰余は他との交叉以外で非分岐である。
交叉は $(1,1)$ の整曲線または異なる二罫線であり、上と同じ境界Picard計算でその剰余が零になる。
すると他方の剰余もその滑らかな部分全体へ延びて零である。
双方が接超平面切断なら一方の補集合 $\mathbb A^3$ を使う。他方の式は、
その接点が補集合内にあれば階数3の二次錐 $xy-z^2$、外なら一次式である。
実際滑らかな二次曲面三次元体の極性により、二接点が互いの接超平面にあるときだけ二次項が消える。
いずれも先のアフィン計算で指数は高々2。全ての極因子を除いた。
\end{proof}
